\loai{Tính tuổi mẫu chất}
\begin{vidu}%[3P7-7.6K][Loai1]
	\nguon{ĐH - 2012} Hạt nhân uranium \Hn{92}{238}{U} sau một chuỗi phân rã, biến đổi thành hạt nhân chì \Hn{82}{206}{Pb}. Trong quá trình đó, chu kì bán rã của \Hn{92}{238}{U} biến đổi thành hạt nhân chì là $4,47.10^9$ năm. Một khối đá được phát hiện có chứa $1,188.10^{20}$ hạt nhân \Hn{92}{238}{U} và $6,239.10^{18}$ hạt nhân \Hn{82}{206}{Pb}. Giả sử khối đá lúc mới hình thành không chứa chì và tất cả lượng chì có mặt trong đó đều là sản phẩm phân rã của \Hn{92}{238}{U}. Tuổi của khối đá khi được phát hiện là
	\choice
	{\True $3,3.10^8$ năm}
	{$6,3.10^9$ năm}
	{$3,5.10^7$ năm}
	{$2,5.10^6$ năm}
	\loigiai{
		\begin{itemize}
			\item Ta có: $N=N_0e^{-\lambda t}$. 
			\item Số hạt U238 bị phân rã bằng số hạt Pb206: $\Delta N=N_0-N=N_0(1-e^{-\lambda t})$ 
			\item $\dfrac{\Delta N}{N}=\dfrac{1-e^{-\lambda t}}{{e^{-\lambda t}}}\Leftrightarrow \dfrac{6,239.10^{18}}{{1,{188.10}^{20}}}=e^{\lambda t}-1\Rightarrow e^{\tfrac{0,693}{T}t}=1+0,053
			\Rightarrow t\approx 3,3.10^8$ năm.
		\end{itemize}
	}
\end{vidu}
\loai{Tuổi trái đất}
\begin{vidu}%[3P7-7.6G][Loai2]
	Cho biết $\Hn{92}{238}{U}$ và $\Hn{92}{235}{U}$ là các chất phóng xạ có chu kì bán rã lần lượt là $T_1 = 4,5.10^9$ năm và $T_2 = 7,13.10^8$ năm. Hiện nay trong quặng uranium thiên nhiên có lẫn $^{238}U$ và $^{235}U$ theo tỉ lệ 160:1. Giả thiết ở thời điểm tạo thành Trái Đất tỉ lệ 1:1. Tuổi của Trái Đất là
	\choice
	{\True 6,2 tỉ năm}
	{5 tỉ năm}
	{5,7 tỉ năm}
	{6,5 tỉ năm}
	\loigiai{
		Ta có: $\dfrac{N_{U238}}{N_{U235}}=\dfrac{N_{0_{U238}}}{{{N}_{0235}}.2^{t\left( \tfrac{1}{T_2}-\tfrac{1}{T_1} \right)}}\Rightarrow 160={{1. 2}^{t\left( \tfrac{1}{7,13.10^8}-\tfrac{1}{4,5.10^{9}} \right)}}\Rightarrow t=6,2.10^{9}$ năm.
	}
\end{vidu}
\loai{So sánh tuổi của 2 mẫu chất}
\begin{vidu}%[3P7-8.5K][Loai3]
	Có hai mẫu chất phóng xạ A và B thuộc cùng đồng vị phóng xạ của một chất có chu kì bán rã $T = 138,2$ ngày và có khối lượng ban đầu như nhau. Tại thời điểm khảo sát, tỉ số hạt nhân hai mẫu chất $\dfrac{N_A}{N_B}=2,72$. Tuổi của mẫu A ít hơn tuổi mẫu B là
	\choice
	{$199,8$ ngày}
	{$189,8$ ngày}
	{\True $199,5$ ngày}
	{$190,4$ ngày}
	\loigiai{
		\begin{itemize}
			\item 	Vì $m_{0A}=m_{0B}$ nên $N_{0A}=N_{0B}$ và  
			$\left\{\begin{aligned}& N_A=N_{0A}.2^{-\tfrac{t_A}{T}}\quad (1) \\
			& N_B=N_{0B}.2^{-\tfrac{t_B}{T}}\quad (2) 
			\end{aligned}\right.$ \\
			\item Lấy $\dfrac{(2)}{(1)}$ ta được: $t_A-t_B= -T\log_2 \dfrac{N_A}{N_B} = 199,5$.
		\end{itemize}
	}
\end{vidu}
\loai{Xạ trị}
\begin{vidu*}%[3P7-7.6G][Loai5]
	Một bệnh nhân phải xạ trị (điều trị bằng đồng vị phóng xạ), dùng tia gamma để diệt tế bào bệnh. Thời gian chiếu xạ lần đầu là 20 phút, cứ sau 1 tháng thì bệnh nhân phải tới bệnh viện để xạ trị. Biết đồng vị phóng xạ đó có chu kì bán rã 4 tháng và vẫn dùng nguồn phóng xạ ban đầu. Hỏi lần chiếu xạ thứ 4 phải có thời gian chiếu xạ là bao lâu để bệnh nhân nhận được lượng tia gamma như lần đầu?
	\choice
	{20 phút}
	{\True 33,6 phút}
	{24,4 phút}
	{40 phút}	
	\loigiai{
		\begin{itemize}
			\item Lượng phóng xạ mà người đó nhận được trong lần đầu: $\Delta N=N_0\pbrace{1-2^{-\tfrac{t_1}{T}}}$.
			\item Với lần xạ trị thứ 4, khoảng thời gian trôi qua là 3 tháng, vậy lượng chất phóng xạ ấy còn lại là: $N=N_02^{-\tfrac{3}{4}}$.
			\item Lượng phóng xạ mà người đó nhận được trong lần thứ tư $\Delta N'=N_02^{-\tfrac{3}{4}}\pbrace{1-2^{-\tfrac{t_2}{T}}}$.
			\item Theo giả thuyết
			\begin{align*}
				\Delta N=\Delta N'\Leftrightarrow N_02^{-\tfrac{3}{4}}\pbrace{1-2^{-\tfrac{t_2}{T}}}=N_0\pbrace{1-2^{-\tfrac{t_1}{T}}}\Leftrightarrow 2^{-\tfrac{3}{4}}\pbrace{1-2^{-\tfrac{t_2}{T}}}=\pbrace{1-2^{-\tfrac{20}{T}}}\Rightarrow t\approx 33\ \text{phút}.
			\end{align*}
		\end{itemize}
	}
\end{vidu*}
